\textit{Syllabus note: This section is provided for enrichment only. Complex-valued methods for mechanics are outside the scope of NSW HSC Mathematics Extension~1 and Extension~2, but they give a useful preview of undergraduate mathematics, physics, and engineering.}

\subsubsection*{1. The complex exponential $e^z$}
In Extension~2, the exponential function is first met for real inputs. If we extend the exponent to a complex number $z=x+iy$, Euler's formula gives
\[
e^{iy}=\cos y+i\sin y.
\]
Hence
\[
e^z=e^{x+iy}=e^x(\cos y+i\sin y).
\]
This is the bridge between oscillation and exponential form: sine-cosine motion can be packaged into one algebraically convenient object.

\subsubsection*{2. Homogeneous differential equations and imaginary roots}
Consider the simple harmonic motion equation
\[
\frac{d^2y}{dx^2}+n^2y=0.
\]
Its characteristic equation is
\[
\lambda^2+n^2=0,
\]
so the roots are $\lambda=\pm ni$. The complex-form solution is
\[
y=C_1e^{inx}+C_2e^{-inx}.
\]
Applying Euler's formula turns this back into the familiar real form
\[
y=A\cos(nx)+B\sin(nx).
\]
So the HSC trigonometric answer is really the real-valued version of a compact exponential calculation.

\subsubsection*{3. Streamlining 2D resistive motion}
In projectile motion with linear resistance $R=-mkv$, the standard Extension~2 approach solves
\[
\ddot{x}=-k\dot{x},
\qquad
\ddot{y}=-g-k\dot{y}
\]
as two separate equations. If we instead define the complex displacement
\[
z=x+iy,
\]
then the complex velocity is $\dot{z}=\dot{x}+i\dot{y}$ and
\[
\ddot{z}=\ddot{x}+i\ddot{y}=-k(\dot{x}+i\dot{y})-ig.
\]
Therefore
\[
\ddot{z}=-k\dot{z}-ig.
\]
This combines the whole 2D system into one linear equation for the complex velocity.

\subsubsection*{Case study: the elegance of the complex plane}
\textbf{Traditional split:}
\[
\dot{x}(t)=u\cos\theta\,e^{-kt},\qquad
\dot{y}(t)=\left(u\sin\theta+\frac{g}{k}\right)e^{-kt}-\frac{g}{k}.
\]

\textbf{Complex method:} Let $W=\dot{z}$. Then
\[
\dot{W}+kW=-ig.
\]
Using the integrating factor $e^{kt}$,
\[
\frac{d}{dt}(We^{kt})=-ig\,e^{kt}
\]
so
\[
W(t)=Ce^{-kt}-\frac{ig}{k}.
\]
Since
\[
W(0)=u(\cos\theta+i\sin\theta)=ue^{i\theta},
\]
we get
\[
C=ue^{i\theta}+\frac{ig}{k}.
\]
Hence the full complex velocity is
\[
\boxed{W(t)=\left(ue^{i\theta}+\frac{ig}{k}\right)e^{-kt}-\frac{ig}{k}}.
\]
Taking the real and imaginary parts recovers the horizontal and vertical velocity components at once.

\subsubsection*{Why this matters}
\begin{itemize}[leftmargin=*]
  \item Complex numbers are not just abstract algebra; they can compress a two-equation mechanics problem into one line.
  \item Imaginary roots in differential equations are the reason sine and cosine appear in oscillation models.
  \item This style of thinking becomes standard in university differential equations, control, waves, and electrical engineering.
\end{itemize}
